\chapter{Cơ sở lý thuyết và khác biệt khi chuyển đổi}
\section{Quy trình huấn luyện}
Phân biệt API mức cao với vòng lặp tường minh PyTorch, logits, softmax và hàm mất mát.

\section{Định dạng ảnh NHWC và NCHW}
Ảnh Keras thường có shape $(N,H,W,C)$; các phép chập PyTorch nhận $(N,C,H,W)$. Phép đổi trục là
\[
x_{\mathrm{NCHW}}=\operatorname{permute}(x_{\mathrm{NHWC}},0,3,1,2).
\]
Với WIKI, phép xoay legacy $\texttt{np.rot90}(k=3,\texttt{axes=(1,2)})$ xảy ra trước khi đổi layout.

\section{CNN kết hợp LSTM}
Thay \texttt{TimeDistributed(CNN)} bằng gộp batch và thời gian $(B,T,C,H,W)\to(BT,C,H,W)$, trích đặc trưng rồi khôi phục $(B,T,F)$ cho LSTM. Tương đương shape không tự chứng minh tương đương trọng số hoặc metric.

\section{Upsampling và phân đoạn}
Đối chiếu \texttt{UpSampling2D} với \texttt{nn.Upsample}, logits theo lớp và mask class-index cho CrossEntropy. Báo riêng total, trainable, non-trainable theo quy ước Keras.

\section{FCN và chuẩn hóa}
Giữ chuẩn hóa legacy $x/127.5-1.0$ khi áp dụng đúng đầu vào. Score-map hồi quy không tự động là detector được đánh giá bằng IoU/mAP.

\section{Chuẩn hóa và giao thức đánh giá}
Nêu chuẩn hóa theo dataset, tách train/validation/test và cảnh báo khi notebook tái sử dụng validation/public-test để in metric mang nhãn Test.